\documentclass[11pt]{article}
\usepackage[a4paper,margin=21mm]{geometry}
\usepackage[no-math]{fontspec}
\setmainfont{Latin Modern Roman}
\newfontfamily\CJKfont{Noto Serif CJK SC}
\XeTeXlinebreaklocale "zh"
\XeTeXlinebreakskip=0pt plus 1pt
\usepackage{amsmath,amssymb,amsthm,amscd,graphicx,color}
\usepackage{xurl}
\usepackage[unicode,hidelinks]{hyperref}
\allowdisplaybreaks
\setlength\emergencystretch{3em}
\numberwithin{equation}{section}

\usepackage{subdepth}

\numberwithin{equation}{section}

\newtheorem{Theorem}{Theorem}[section]
\newtheorem*{Theorem*}{Theorem}
\newtheorem{Corollary}[Theorem]{Corollary}
\newtheorem{Lemma}[Theorem]{Lemma}
\newtheorem{Proposition}[Theorem]{Proposition}
 { \theoremstyle{definition}
\newtheorem{Definition}[Theorem]{Definition}
\newtheorem{Note}[Theorem]{Note}
\newtheorem{Example}[Theorem]{Example}
\newtheorem{Remark}[Theorem]{Remark} }

\def\C {\mathbb{C}}
\def\E {\mathbb{E}}
\def\N {\mathbb{N}}
\def\R {\mathbb{R}}
\def\Var {\mathrm{Var}}
\def\Cov {\mathrm{Cov}}
\def\dd {\hskip1pt\mathrm{d}}
\def\tr {\mathop{\mathrm{tr}}}
\def\cid{\stackrel{\mathcal{L}}\to}
\def\cip{\stackrel{\mathcal{P}}\to}
\def\diag{\mathop{\mathrm{diag}}}
\def\Or{\mathop{\mathrm{Or}}}
\def\mcalS{\mathcal{S}}

\newcommand{\wtilde}{\widetilde}
\newcommand{\eqdist}{\overset{\mathcal{L}}{=}}
\newcommand{\stge}{\overset{\mathcal{L}}{\ge}}
\newcommand{\stle}{\overset{\mathcal{L}}{\le}}
\newcommand{\bfone}{\boldsymbol{1}}
\newcommand{\bzero}{\boldsymbol{0}}
\newcommand{\Proof}{\noindent\textit{Proof.}\ \ }
\newcommand{\sym}{{{\hskip-1pt}\mathop{\mathrm{sym}}}}

\begin{document}
\begin{center}\Large Explicit all-dimension remainder bounds for matrix Bingham normalization with symmetric parameters and sub-square-root Frobenius growth\end{center}
\noindent\textbf{Stable ID:} 1051\quad\textbf{Author:} Armine Bagyan; Donald Richards\par
\noindent\textbf{Work:} Complete Asymptotic Expansions for the Normalizing Constants of High-Dimensional Matrix Bingham and Matrix Langevin Distributions\par
\noindent\textbf{Source:} \url{https://sigma-journal.com/2024/094/}\par
\noindent\textbf{Version:} Official SIGMA 20 (2024) article 094, DOI 10.3842/SIGMA.2024.094; journal PDF and native TeX acquired 2026-09-30, exact bytes pinned by SHA256\par
\noindent\textbf{Rights:} CC BY 4.0. Authors retain copyright. \url{https://creativecommons.org/licenses/by/4.0/}. Reuse and typesetting adaptation; original language and mathematical content preserved.\par
\noindent\textbf{Difficulty (prior selection preserved):} {\CJKfont H2: The proof controls a matrix hypergeometric tail through zonal-polynomial integrals over the orthogonal group, combining principal-minor bounds with marginal Haar moments. Partition-wise shifted-factorial estimates then give constants uniform across the summation, and a telescoping factorial-tail argument produces the explicit closed exponential bound. Both the series and closed-form estimates belong to one proof unit for d at least p at least one and Frobenius growth exponent 0 <= r < 1.}\par
\medskip\noindent\textbf{Editorial boundary note (not author text):} Complete original Theorem2.1 for all d at least p at least1, symmetric parameters, every Frobenius growth condition, m at least2, both the infinite-series and exact closed exponential bounds, and the stated remainder rate. All partition/zonal/Frobenius/Haar definitions and current supporting proofs are retained: full factorial-product and shifted-partition inequalities, nonnegative monomial/absolute-eigenvalue domination, principal-minor Hadamard bound, Hölder factorization, marginal Haar moments and full zonal estimate, followed by the complete partition sum and telescoping factorial-tail proof with both Cauchy–Schwarz estimates and exact constants. James integral/normalization formulas and precise earlier one-part estimates Bagyan–Richards LemmaA.2/equationA.10 remain displayed, precisely cited prerequisites; the current new all-partition and all-dimension constructions are fully present. The infinite-series and closed-form estimates count as one primary remainder problem, positioned as a main article outcome. Lower-bound results, Langevin outcomes, unimodality comments, numerical examples and applications are excluded, not additional counts. Original equations, including literal apparent reciprocal/bar quirks, remain unchanged without assistant verification, proof generation or formula repair. No image-only mathematical content.\par
\medskip\hrule\medskip\noindent\textbf{Author text begins below.}\par
\setcounter{section}{1}
% Original source lines 84--91
For positive integers $d$ and $p$ such that $d \ge p$ let $\R^{d \times p}$ denote the set of $d \times p$ real matrices, $I_p$ be the identity matrix of order $p$, and \smash{$V_{d,p} = \bigl\{x \in \R^{d \times p} \mid x'x = I_p\bigr\}$} denote the \textit{Stiefel manifold}. The manifold $V_{d,p}$ is endowed with a unique probability measure, denoted by $\dd x$, that is invariant under the change of variables $x \mapsto HxK$, where $x \in V_{d,p}$, and $H$ and $K$ are orthogonal matrices of orders $d \times d$ and $p \times p$, respectively (see Chikuse \cite[p.~16]{Chikuse}, Muirhead \cite[p.~70]{Muirhead}).

In this article, we first consider the matrix Bingham probability distributions on $V_{d,p}$. For non-zero symmetric matrices \smash{$A \in \R^{p \times p}$} and \smash{$\Sigma \in \R^{d \times d}$}, a random matrix $X \in V_{d,p}$ is said to have a {\it $($full$)$ matrix Bingham distribution with parameters $A$ and $\Sigma$} if the probability density function of $X$, with respect to the invariant measure $\dd x$, is
\begin{equation}
\label{eq_mBingham_pdf}
\phi(x;A,\Sigma) = [\Phi_{d,p}(A,\Sigma)]^{-1} \exp\bigl(\tr Ax'\Sigma x\bigr), \qquad x \in V_{d,p},
\end{equation}
where $\Phi_{d,p}(A,\Sigma)$ is the normalizing constant. The matrix Bingham family is now ubiquitous in the statistical analysis of directional data and shape analysis, and the extant literature now provides many properties of the family and its applications in fields including astronomy, biology, computer vision, medicine, meteorology, physics and psychology; see, e.g., \cite{BinghamCR,Chikuse,Dryden,JuppMardia2,KhatriMardia,KumeWood,KumePrestonWood,MardiaJupp,Prentice}.


% Original source lines 113--165
\section[Complete asymptotics for the normalizing constant of the high-dimensional matrix Bingham distribution]{Complete asymptotics for the normalizing constant\\ of the high-dimensional matrix Bingham distribution}
\label{sec_matrix_bingham}

A \textit{partition} $\kappa = (\kappa_1,\ldots,\kappa_d)$ is a vector of integers such that $\kappa_1 \ge \cdots \ge \kappa_d \ge 0$. The non-zero~$\kappa_j$~are called the \textit{parts} of $\kappa$. The \textit{length} of $\kappa$, denoted by $\ell(\kappa)$, is the number of non-zero~$\kappa_j$,~$j=1,\ldots,d$, and the \textit{weight} of $\kappa$ is $|\kappa| = \kappa_1+\cdots+\kappa_d$.

Define the \textit{shifted factorial}, $(a)_k = a(a+1)(a+2) \cdots (a+k-1)$, where $a \in \mathbb{R}$ and ${k=0,1,2,\ldots}$; in particular, $(a)_0 \equiv 1$ since the defining product is empty. For $a \in \mathbb{R}$ and a partition ${\kappa = (\kappa_1,\ldots,\kappa_d)}$, the {\it partitional shifted factorial} is %defined as
\begin{equation}
\label{eq_partit_fact}
(a)_\kappa = \prod_{i=1}^d \left(a-\frac{1}{2}(i-1)\right)_{\kappa_i}.
\end{equation}

Assume that \smash{$\R^{d \times d}_\sym$} denote the space of real symmetric $d \times d$ matrices. For a partition $\kappa = (\kappa_1,\ldots,\kappa_d)$ and \smash{$\Sigma \in \R^{d \times d}_\sym$}, the \textit{zonal polynomial} $C_\kappa(\Sigma)$ is an eigenfunction of the Laplace--Beltrami operator in the eigenvalues of $\Sigma$ (see \cite{James68}, \cite[Chapter 7]{Muirhead} and \cite{Richards10}). Further, $C_\kappa(\Sigma)$~depends only on the eigenvalues of $\Sigma$, and is symmetric in those eigenvalues.

Tables of the polynomials $C_\kappa(\Sigma)$ for $|\kappa| \le 12$ were provided by James \cite[pp.~498--499]{James64} and Parkhurst and James \cite{ParkhurstJames}. Jiu and Koutschan \cite{JiuKoutschan} have provided software to calculate rapidly the zonal polynomials to high values of $|\kappa|$, so the asymptotic expansions obtained in this article can be calculated for practical values of $m$.

Let $X \in V_{d,p}$ be a random matrix having the matrix Bingham distribution with the density function $\phi(x;A,\Sigma)$ as in \eqref{eq_mBingham_pdf}. It is well known that the normalizing constant $\Phi_{d,p}(A,\Sigma)$ can be expressed as a generalized hypergeometric function of two matrix arguments, denoted by~${}_0F_0(A,\Sigma)$, or by the zonal polynomial expansion of that function,
\begin{equation}
\label{eq_normaliz_const_B}
\Phi_{d,p}(A,\Sigma) = {}_0F_0(A,\Sigma) = \sum_{k=0}^\infty \frac{1}{k!} \sum_{|\kappa|=k} \frac{C_\kappa(A) C_\kappa(\Sigma)}{C_\kappa(I_d)},
\end{equation}
where the inner sum is over all partitions $\kappa$ such that $\ell(\kappa) \le p$. This formula for $\Phi_{d,p}(A,\Sigma)$ in terms of ${}_0F_0(A,\Sigma)$ dates to De Waal \cite{DeWaal}; cf.\ Bingham \cite{Bingham}, Chikuse \cite[p.~108]{Chikuse}. By applying the methods of \cite[Theorem~6.3]{GrossRichards}, we find that the series \eqref{eq_normaliz_const_B} converges absolutely for all $(A,\Sigma)$ and uniformly on compact regions.

For the special case in which $A = I_p$, on applying to \eqref{eq_normaliz_const_B} the identity \eqref{eq_kappa_invar}, written in the form $C_\kappa(I_p)/C_\kappa(I_d) = (p/2)_\kappa/(d/2)_\kappa$, we obtain the normalizing constant in terms of a confluent hypergeometric function of matrix argument,
%\label{eq_normaliz_const_hgf}
\[
\Phi_{d,p}(I_p,\Sigma) = \sum_{k=0}^\infty \frac{1}{k!} \sum_{|\kappa|=k} \frac{(p/2)_\kappa}{(d/2)_\kappa} C_\kappa(\Sigma) = {}_1F_1\left(\frac12 p;\frac12 d;\Sigma\right),
\]
as shown by various authors, e.g., Bingham \cite{Bingham}, Chikuse \cite[p.~33]{Chikuse} and Muirhead \cite[p.~288]{Muirhead}.

Similar to \cite[equation~(2.4)]{BagyanRichards}, when $d$ is large, our asymptotic approximation to $\Phi_{d,p}(A,\Sigma)$ is the sum of the terms up to degree $m-1$ of the series \eqref{eq_normaliz_const_B}, i.e.,
\begin{equation}
\label{eq_trunc_expan}
\Phi_{d,p}(A,\Sigma) \approx \sum_{k=0}^{m-1} \frac{1}{k!} \sum_{|\kappa|=k} \frac{C_\kappa(A) C_\kappa(\Sigma)}{C_\kappa(I_d)},
\end{equation}
$m \ge 2$, and now we study the remainder series
\begin{equation}
\label{eq_remainder_term}
\Phi_{d,p;m}(A,\Sigma) = \sum_{k=m}^\infty \frac{1}{k!} \sum_{|\kappa|=k} \frac{C_\kappa(A) C_\kappa(\Sigma)}{C_\kappa(I_d)}.
\end{equation}

In the sequel, we will encounter the constants \smash{$\alpha_p = (2\pi)^{-(p-1)/4p} p^{1/4p}$} and $\gamma_1 = \bigl(\sqrt{3}+1\bigr)/2$, arising in Lemmas \ref{lemma_factorial_ineq} and \ref{lemma_ineqs}, respectively. Also, if $a_1,\ldots,a_p$ are the eigenvalues of $A$, then we~define $A_+ = \diag(|a_1|,\ldots,|a_p|)$, the diagonal matrix with diagonal entries $|a_1|,\ldots,|a_p|$. In~the following result, the proof of which is given in Section \ref{sec_proofs_Phi_bounds}, we extend a result in~\cite{BagyanRichards} by~deriving an upper bound for $\Phi_{d,p;m}(A,\Sigma)$ in terms of the Frobenius norm, \smash{$\|\Sigma\| = \bigl[\tr\bigl(\Sigma^2\bigr)\bigr]^{1/2}$}.

\begin{Theorem}
\label{thm_Phi_m_bound_Bingham}
Let $A \in \R^{p \times p}_\sym$, \smash{$\Sigma \in \R^{d \times d}_\sym$}, and suppose that \smash{$\|\Sigma\| \le \gamma_0 d^{r/2}$}, where $\gamma_0 > 0$ and $0 \le r < 1$. Then for all $m \ge 2$,
\begin{align}
|\Phi_{d,p;m}(A,\Sigma)| \le{}& \alpha_p \sum_{k=m}^\infty \frac{\bigl[\gamma_0 \gamma_1 p^{1/2} (\tr A_+) d^{-(1-r)/2}\bigr]^k}{(k!)^{1/2}} \label{eq_Phi_m_bound} \\
\le{}& \alpha_p (4 {\rm e}/\pi)^{1/4} ({\rm e}/m)^{(m/2)-(1/4)}
\bigl[\gamma_0 \gamma_1 p^{1/2} (\tr A_+) d^{-(1-r)/2}\bigr]^m \nonumber \\
& \times \exp\bigl(c_m \bigl[\gamma_0 \gamma_1 p^{1/2} (\tr A_+) d^{-(1-r)/2}\bigr]^2/2\bigr), \label{eq_Phi_m_boundw}
\end{align}
where $c_m = \bigl(1+m^{-1}\bigr)^{-m} {\rm e}$. In particular, $\Phi_{d,p;m}(A,\Sigma) = O\bigl(d^{-(1-r)m/2}\bigr)$ as $d \to \infty$.
\end{Theorem}

\setcounter{section}{4}
% Original source lines 324--599
\section{Proofs}\label{sec_proofs}

\subsection{Properties of the zonal polynomials}\label{sec_zonal_polys}

James \cite{James73} obtained for the zonal polynomials $C_\kappa(\Sigma)$ a crucial integral formula, as we now explain. Let \smash{$\Or(d) = \bigl\{H \in \R^{d \times d}\mid H'H = I_d\bigr\}$} be the group of orthogonal matrices in \smash{$\R^{d \times d}$}. For~${H \in \Or(d)}$, denote by $\dd H$ the Haar measure on $\Or(d)$, normalized to have total volume $1$. Also let ${\det}_j(\Sigma)$ be the $j$th principal minor of $\Sigma$, $j=1,\ldots,d$. For each partition $\kappa = (\kappa_1,\ldots,\kappa_d)$,
\begin{gather}
\label{eq_zonal}
C_\kappa(\Sigma) = C_\kappa(I_d) \int_{\Or(d)} \prod_{j=1}^{\ell(\kappa)} \bigl[{\det}_j\bigl(H' \Sigma H\bigr)\bigr]^{\kappa_j-\kappa_{j+1}} \dd H,
\end{gather}
where $\kappa_{d+1} \equiv 0$ and
\begin{equation}
\label{eq_zonal_identity}
C_\kappa(I_d) = 2^{2 |\kappa|} |\kappa|! (d/2)_\kappa \frac{\prod\limits_{1 \le i < j \le \ell(\kappa)} (2\kappa_i - 2\kappa_j - i + j)}{\prod\limits_{i=1}^{\ell(\kappa)} (2\kappa_i + \ell(\kappa) - i)!}.
\end{equation}
The formula \eqref{eq_zonal} was first given by James \cite[equation~(10.4)]{James73}, further applications of it were given in \cite{Kushner} and \cite{Richards85}, and the formula was extended to a more general setting in \cite[Theorem~4.8]{GrossRichards}.

By \eqref{eq_zonal}, $C_\kappa(\Sigma)$ is homogeneous of degree $|\kappa|$; moreover $C_\kappa(\Sigma)$ depends only on the $\ell(\kappa)$ non-zero parts of $\kappa$, so that $C_\kappa(\Sigma)$ is unchanged if zeros are appended to $\kappa$. Since any symmetric matrix can be diagonalized by an orthogonal transformation, then it also follows from \eqref{eq_zonal} that~$C_\kappa(\Sigma)$ is a symmetric function of the eigenvalues of $\Sigma$.

The importance of the normalization \eqref{eq_zonal_identity} is that the zonal polynomials satisfy the identity,
\begin{equation}
\label{eq_trace_zonal}
(\tr \Sigma)^k = \sum_{|\kappa|=k} C_\kappa(\Sigma),
\end{equation}
for all $k=0,1,2,\ldots$; see \cite[p.~479, equation~(21)]{James64}, \cite[p.~228, equation~(iii)]{Muirhead}, \cite[equation~(35.4.6)]{Richards10}. Also noteworthy is the special case of \eqref{eq_zonal_identity} in which $\kappa = (k)$, a partition with one part
\begin{equation}
\label{eq_zonal_identity_1part}
C_{(k)}(I_d) = \frac{(d/2)_k}{(1/2)_k}.
\end{equation}

The (seemingly formidable) expression for $C_\kappa(I_d)$ in \eqref{eq_zonal_identity} arises from the representation theory of $\mathfrak{S}_k$, the group of permutations on $k$ symbols, and we comment on the connection as follows. It is a classical result that the collection of irreducible representations of $\mathfrak{S}_k$ is in one-to-one correspondence with the set of partitions of $k$ (see Ledermann \cite[Chapter 4]{Ledermann}). For a~given partition $\kappa = (\kappa_1,\ldots,\kappa_d)$ we use the standard notation $\chi_{[2\kappa]}(1)$ for the dimension of the irreducible representation, corresponding to the partition $2\kappa = (2\kappa_1,\ldots,2\kappa_d)$, of the symmetric group $\mathfrak{S}_{2|\kappa|}$ on $2|\kappa|$ symbols. James \cite[pp.~478--479, equations~(18) and~(20)]{James64} proved that
\begin{equation}
\label{eq_zonal_identity_rep}
C_\kappa(I_d) = \frac{2^{2|\kappa|} |\kappa|! (d/2)_\kappa}{(2|\kappa|)!} \chi_{[2\kappa]}(1) = \frac{(d/2)_\kappa}{(1/2)_{|\kappa|}} \chi_{[2\kappa]}(1),
\end{equation}
and then \eqref{eq_zonal_identity} is derived from \eqref{eq_zonal_identity_rep} by applying Frobenius' famous formula for $\chi_{[2\kappa]}(1)$ (Ledermann \cite[p.~123]{Ledermann}).

Another consequence of \eqref{eq_zonal_identity_rep} is that
\begin{equation}
\label{eq_chi_kappa}
\frac{C_\kappa(I_d)}{(d/2)_\kappa} = \frac{\chi_{[2\kappa]}(1)}{(1/2)_{|\kappa|}}.
\end{equation}
Since the dimension of an irreducible representation does not change if zeros are appended to the corresponding partition, then it follows that, for any partition $\kappa$ such that $\ell(\kappa) \le p$, the right-hand side of \eqref{eq_chi_kappa} remains unchanged if $d$ is replaced by $p$; hence
\begin{equation}
\label{eq_kappa_invar}
\frac{C_\kappa(I_d)}{(d/2)_\kappa} = \frac{C_\kappa(I_p)}{(p/2)_\kappa}.
\end{equation}
We remark that the identity \eqref{eq_kappa_invar} can be verified directly, albeit laboriously, from \eqref{eq_zonal_identity}. Further, \eqref{eq_kappa_invar} has been noted explicitly in the literature, e.g., by Chikuse \cite[p.~30]{Chikuse} and by Edelman et al. \cite[p.~259]{Edelman}. Also, if $\ell(\kappa) = 1$, then \eqref{eq_kappa_invar} reduces to \eqref{eq_zonal_identity_1part}.

The zonal polynomials also have a combinatorial formulation. Recall that the set of all partitions is endowed with the \textit{lexicographic ordering}: For partitions $\kappa = (\kappa_1,\ldots,\kappa_d)$ and ${\lambda = (\lambda_1,\ldots,\lambda_d)}$ such that $|\kappa| = |\lambda|$, we say that $\lambda$ \textit{is less than} $\kappa$, written $\lambda < \kappa$, if there exists $i \in \{1,\ldots,d-1\}$ such that $(\lambda_1,\ldots,\lambda_{i-1}) = (\kappa_1,\ldots,\kappa_{i-1})$ and $\lambda_i < \kappa_i$.

For \smash{$\Sigma \in \R^{d \times d}_\sym$}, denote by $\sigma_1,\ldots,\sigma_d$ the eigenvalues of $\Sigma$. For each partition $\kappa = (\kappa_1,\ldots,\kappa_\ell)$ of length $\ell$, the {\it monomial symmetric function of $\Sigma$ corresponding to $\kappa$} is
\begin{equation}
\label{eq_msf}
M_\kappa(\Sigma) = \sum_{(j_1,\ldots,j_\ell)} \sigma_{j_1}^{\kappa_1} \cdots \sigma_{j_\ell}^{\kappa_\ell},
\end{equation}
where the sum is taken over all distinct permutations $(j_1,\ldots,j_\ell)$ of $\ell$ integers drawn from the set $\{1,\ldots,d\}$. James \cite{James68} (see also Muirhead \cite[Section 7.2]{Muirhead}) proved that
\begin{equation}
\label{eq_zonal_msf}
C_\kappa(\Sigma) = \sum_{\lambda \le \kappa} c_{\kappa,\lambda} M_\lambda(\Sigma),
\end{equation}
where the constants $c_{\kappa,\lambda}$ are \textit{nonnegative} and the summation is over all partitions $\lambda$ such that~${\lambda \le \kappa}$, i.e., $\lambda$ is less than or equal to $\kappa$ in the lexicographic ordering.


\subsection[Inequalities for the partitional shifted factorials and the zonal polynomials]{Inequalities for the partitional shifted factorials\\ and the zonal polynomials}
\label{sec_ineqs}

\begin{Lemma}
\label{lemma_factorial_ineq}
Let $a_1,\ldots,a_p$ be nonnegative integers such that $a_1+\cdots+a_p\ge 2$, and define
\begin{equation}
\label{eq_alpha_p}
\alpha_p = (2\pi)^{-(p-1)/4p} p^{1/4p}.
\end{equation}
Then
\begin{equation}
\label{eq_prod_factorial_bound}
\prod_{i=1}^p (p a_i)! \le \alpha_p^{2p} p^{p(a_1+\cdots+a_p)} \bigl[(a_1+\cdots+a_p)!\bigr]^p.
\end{equation}
\end{Lemma}

\begin{proof}
Since the multinomial coefficient, $(a_1+\cdots+a_p)/a_1! \cdots a_p!$, is a positive integer, then
\begin{equation}
\label{eq_multinomial}
\prod_{i=1}^p a_i! \le (a_1 + \cdots + a_p)!.
\end{equation}
It is well known that the classical gamma function is increasing on the interval $[2,\infty)$; therefore,
\smash{$\Gamma\bigl(a+(i-1)p^{-1}\bigr) \le \Gamma\bigl(a+(p-1)p^{-1}\bigr) = \Gamma\bigl(a + 1 - p^{-1}\bigr)$}
for $a \ge 2$ and all $i=1,\ldots,p$. Applying Gauss' multiplication formula for the gamma function, we obtain
\begin{align}
\Gamma(pa) &= (2\pi)^{-(p-1)/2} p^{pa-(1/2)} \prod_{i=1}^p \Gamma\bigl(a + (i-1)p^{-1}\bigr) \nonumber \\
&\le (2\pi)^{-(p-1)/2} p^{pa-(1/2)} \bigl[\Gamma\bigl(a + 1 - p^{-1}\bigr)\bigr]^p,\label{eq_Gauss_mult_1}
\end{align}
where the latter inequality holds for $a \ge 2$.

Setting $a = a_1+\cdots+a_p + p^{-1}$ in \eqref{eq_Gauss_mult_1}, we obtain
\begin{align}
\label{eq_Gauss_mult_2}
(p (a_1+\cdots+a_p))! &= \Gamma(p (a_1+\cdots+a_p) + 1) \nonumber \\
&\le (2\pi)^{-(p-1)/2} p^{p (a_1+\cdots+a_p) + (1/2)} [\Gamma(a_1+\cdots+a_p+1)]^p \nonumber \\
&= (2\pi)^{-(p-1)/2} p^{p (a_1+\cdots+a_p) + (1/2)} [(a_1+\cdots+a_p)!]^p.
\end{align}
Applying \eqref{eq_Gauss_mult_2} to \eqref{eq_multinomial}, we find that
\begin{align*}
\prod_{i=1}^p (p a_i)! &\le \bigl(p (a_1+\cdots+a_p)\bigr)! \le (2\pi)^{-(p-1)/2} p^{p (a_1+\cdots+a_p) + (1/2)} \bigl((a_1+\cdots+a_p)!\bigr)^p \\
&\equiv \alpha_p^{2p} p^{p(a_1+\cdots+a_p)} [(a_1+\cdots+a_p)!]^p.
\end{align*}
This completes the proof.
\end{proof}



\begin{Lemma}
\label{lemma_ineqs}
Let $d \ge p$ and let $\kappa = (\kappa_1,\ldots,\kappa_p)$ be a partition of length $\ell(\kappa) \le p$ and weight~${|\kappa| \ge 2}$. Then
\begin{equation}
\label{eq_d_kappa_ineq}
(2p)^{-|\kappa|} d^{|\kappa|} \le (d/2)_\kappa \le 2^{-|\kappa|} (d + |\kappa|)^{|\kappa|}.
\end{equation}
Also, with $\gamma_1 = \bigl(\sqrt{3}+1\bigr)/2 \simeq 1.366025$, we have
\begin{equation}
\label{eq_d_kappa_ratio_ineq3}
\prod_{i=1}^p \frac{\bigl(d^{1/2}/2\bigr)_{p \kappa_i}}{(d/2)_{p \kappa_i}} \le \bigl[\alpha_p (|\kappa|!)^{1/2} \gamma_1^{|\kappa|} p^{|\kappa|/2} d^{-|\kappa|/2}\bigr]^p.
\end{equation}
\end{Lemma}

\begin{proof}
By \eqref{eq_partit_fact}, we have
\begin{equation}
\label{eq_d_kappa_ineq_2}
(d/2)_\kappa = \prod_{j=1}^p \left(\frac12 d - \frac12(j-1)\right)_{\kappa_j}
= \prod_{j=1}^p \prod_{i=1}^{\kappa_j} \left(\frac12d-\frac12(j-1)+i-1\right).
\end{equation}
For all $i \ge 1$, $j \le p$, and $d \ge p$, we also have
\begin{align*}
\frac12d-\frac12(j-1)+i-1 &\ge \frac12d-\frac12(p-1) = \frac12 \left(1-\frac{p-1}{d}\right) d
\ge \frac12 \left(1-\frac{p-1}{p}\right) d = (2p)^{-1}d.
\end{align*}
Therefore, by \eqref{eq_d_kappa_ineq_2},
\[
(d/2)_\kappa \ge \prod_{j=1}^p \prod_{i=1}^{\kappa_j} \bigl[(2p)^{-1} d\bigr] = (2p)^{|\kappa|} d^{-|\kappa|}.
\]
This establishes the lower bound in \eqref{eq_d_kappa_ineq}.

Next, we apply to the product in \eqref{eq_d_kappa_ineq_2} the arithmetic-geometric mean inequality. Then we obtain
\begin{align*}
(d/2)_\kappa &\le \Bigg(\frac{1}{|\kappa|} \sum_{j=1}^p \sum_{i=1}^{\kappa_j} \left(\frac12d-\frac12(j-1)+i-1\right)\Bigg)^{|\kappa|} \\
&= \Bigg(\frac{1}{2 |\kappa|} \Bigg(d |\kappa| - \sum_{j=1}^p j \kappa_j + \sum_{j=1}^p \kappa_j^2\Bigg)\Bigg)^{|\kappa|} \le \Bigg(\frac{1}{2 |\kappa|} \Bigg(d |\kappa| + \Bigg(\sum_{j=1}^p \kappa_j\Bigg)^2\Bigg)\Bigg)^{|\kappa|} \\
&= 2^{-|\kappa|} (d + |\kappa|)^{|\kappa|}.
\end{align*}
This proves the upper bound in \eqref{eq_d_kappa_ineq}.

Further, it is known from \cite[Lemma A.2]{BagyanRichards} that, for all $r \ge 1$,
\smash{$\frac{(d^{1/2}/2)_r}{(d/2)_r} \le [(r-1)!]^{1/2} \gamma_1^{r-1} d^{-r/2}$},
and since $\gamma_1 > 1$, then we also have
\[
\frac{\bigl(d^{1/2}/2\bigr)_r}{(d/2)_r} \le [(r-1)!]^{1/2} \gamma_1^{r-1} d^{-r/2} \cdot \gamma_1 r^{1/2} = (r!)^{1/2} \gamma_1^r d^{-r/2}.
\]
Since the latter inequality is valid for $r = 0$, then we have shown that, for all $r \ge 0$,
\begin{equation}
\label{eq_d_kappa_ratio_ineq2}
\frac{\bigl(d^{1/2}/2\bigr)_r}{(d/2)_r} \le (r!)^{1/2} \gamma_1^r d^{-r/2}.
\end{equation}

Now we apply \eqref{eq_d_kappa_ratio_ineq2} to obtain
\begin{equation}
\label{eq_d_kappa_ratio_ineq}
\prod_{i=1}^p \frac{\bigl(d^{1/2}/2\bigr)_{p \kappa_i}}{(d/2)_{p \kappa_i}} \le \prod_{i=1}^p [(p \kappa_i)!]^{1/2} \gamma_1^{p \kappa_i} d^{-p \kappa_i/2}
= \gamma_1^{p |\kappa|} d^{-p |\kappa|/2} \prod_{i=1}^p [(p \kappa_i)!]^{1/2}.
\end{equation}
By \eqref{eq_prod_factorial_bound} in Lemma \ref{lemma_factorial_ineq} with $a_i = \kappa_i$, $i=1,\ldots,p$, we also have
\[
\prod_{i=1}^p (p \kappa_i)! \le \alpha_p^{2p} p^{p |\kappa|} (|\kappa|!)^p,
\]
and by substituting this bound into \eqref{eq_d_kappa_ratio_ineq}, we obtain \eqref{eq_d_kappa_ratio_ineq3}.
\end{proof}

For \smash{$\Sigma \in \R^{d \times d}_\sym$}, let $\sigma_1,\ldots,\sigma_d$ denote the eigenvalues of $\Sigma$ and define the diagonal matrix $\Sigma_+ = \diag(|\sigma_1|,\ldots,|\sigma_d|)$. The following result is a consequence of \eqref{eq_zonal_msf}, the combinatorial formula for the zonal polynomials.

\begin{Lemma}
\label{lem_zonal_psd}
Let $\kappa$ be a partition and \smash{$\Sigma \in \R^{d \times d}_\sym$}. Then $|C_\kappa(\Sigma)| \le C_\kappa(\Sigma_+)$.
\end{Lemma}

\begin{proof}
By \eqref{eq_zonal_msf}, the nonnegativity of the constants $c_{\kappa,\lambda}$, and the triangle inequality, we have
\[
|C_\kappa(\Sigma)| \le \boldsymbol{\bigg|}\sum_{\lambda \le \kappa} c_{\kappa,\lambda} M_\lambda(\Sigma)\boldsymbol{\bigg|}
\le \sum_{\lambda \le \kappa} c_{\kappa,\lambda} |M_\lambda(\Sigma)|.
\]
By \eqref{eq_msf}, for all partitions $\lambda$,
\[
|M_\lambda(\Sigma)| = \boldsymbol{\Big|}\sum \sigma_{j_1}^{\lambda_1} \cdots \sigma_{j_\ell}^{\lambda_\ell}\boldsymbol{\Big|}
\le \sum |\sigma_{j_1}|^{\lambda_1} \cdots |\sigma_{j_\ell}|^{\lambda_\ell}
= M_\lambda(\Sigma_+).
\]
Therefore, \[
|C_\kappa(\Sigma)| \le \sum_{\lambda \le \kappa} c_{\kappa,\lambda} M_\lambda(\Sigma_+) =
C_\kappa(\Sigma_+).
\]
The proof now is complete.
\end{proof}



The following result will be crucial for deriving bounds on the tails of the asymptotic series for the normalizing constants of the matrix Bingham and matrix Langevin distributions.

\begin{Proposition}
\label{prop_zonal_bound}
Let $\kappa$ be a partition such that $\ell(\kappa) \le p$, where $p \le d$, and let \smash{$\Sigma \in \R^{d \times d}_\sym$}. Then
\begin{gather}
\label{eq_zonal_bound}
|C_\kappa(\Sigma)| \le \Bigg(\prod_{i=1}^p \frac{\bigl(d^{1/2}/2\bigr)_{p \kappa_i}}{(d/2)_{p \kappa_i}}\Bigg)^{1/p} C_\kappa(I_d) \|\Sigma\|^{|\kappa|}.
\end{gather}
\end{Proposition}

\begin{proof}
By Lemma \ref{lem_zonal_psd}, $|C_\kappa(\Sigma)| \le C_\kappa(\Sigma_+)$. Since $\|\Sigma\| = \|\Sigma_+\|$, then it suffices to assume, without loss of generality, that $\Sigma = \Sigma_+$, i.e., that $\Sigma$ is positive semidefinite and diagonal, with diagonal entries $\sigma_1,\ldots,\sigma_d$.

Since $\Sigma$ is positive semidefinite, then so is every principal submatrix of $H'\Sigma H$, for all ${H \in \Or(d)}$. Therefore, ${\det}_j\bigl(H'\Sigma H\bigr) \ge 0$ for all $j=1,\ldots,d$.

For $\ell(\kappa) \le p$, it follows from \eqref{eq_zonal} that
\begin{equation}
\label{eq_zonal_ineq1}
C_\kappa(\Sigma) = C_\kappa(I_d) \int_{\Or(d)} \prod_{j=1}^p [{\det}_j\bigl(H'\Sigma H\bigr)]^{\kappa_j-\kappa_{j+1}} \dd H.
\end{equation}
Denote by $\bigl(H'\Sigma H\bigr)_{ii}$ the $i$th diagonal entry of $H'\Sigma H$, $i=1,\ldots,d$. By Hadamard's inequality for the principal minors of a positive semidefinite matrix \cite[p.~505]{Horn},
\[
{\det}_j\bigl(H'\Sigma H\bigr) \le \prod_{i=1}^j \bigl(H'\Sigma H\bigr)_{ii}
\]
for all $j=1,\ldots,d$; therefore,
\begin{align*}
\prod_{j=1}^p \bigl[{\det}_j\bigl(H'\Sigma H\bigr)\bigr]^{\kappa_j-\kappa_{j+1}} &\le \prod_{j=1}^p \prod_{i=1}^j \bigl[\bigl(H'\Sigma H\bigr)_{ii}\bigr]^{\kappa_j-\kappa_{j+1}} \\
&= \prod_{i=1}^p \prod_{j=i}^p \bigl[\bigl(H'\Sigma H\bigr)_{ii}\bigr]^{\kappa_j-\kappa_{j+1}}
= \prod_{i=1}^p \bigl[\bigl(H'\Sigma H\bigr)_{ii}\bigr]^{\kappa_i}.
\end{align*}
Inserting this bound into \eqref{eq_zonal_ineq1} and then applying H\"older's inequality, we obtain
\begin{align}
C_\kappa(\Sigma) &\le C_\kappa(I_d) \int_{\Or(d)} \prod_{i=1}^p \bigl[\bigl(H'\Sigma H\bigr)_{ii}\bigr]^{\kappa_i} \dd H \nonumber \\
&\le C_\kappa(I_d) \Bigg(\prod_{i=1}^p \int_{\Or(d)} \bigl[\bigl(H'\Sigma H\bigr)_{ii}\bigr]^{p \kappa_i} \dd H\Bigg)^{1/p}.\label{eq_zonal_Holder}
\end{align}

Denote by $h_{i,j}$ the $(i,j)$th entry of $H$, then $h_i = (h_{1,i},\ldots,h_{d,i})'$ is the $i$th column of $H$, $i=1,\ldots,d$. Since $\Sigma = \diag(\sigma_1,\ldots,\sigma_d)$, then for all $i=1,\ldots,d$,
\[
\bigl(H'\Sigma H\bigr)_{ii} = \sum_{l=1}^d h_{l,i}^2 \sigma_l = h_i'\Sigma h_i.
\]
Now regard $H$ as a random matrix having the uniform distribution (Haar measure) on $Or(d)$. Since the Haar measure is orthogonally invariant, and therefore invariant under permutation of the columns of $H$, it follows that the marginal distribution of $h_i$ does not depend on $i$. Denoting equality in distribution by \smash{$\eqdist$}, we obtain
\[
\bigl(H'\Sigma H\bigr)_{ii} = h_i'\Sigma h_i \eqdist h_1'\Sigma h_1 = \bigl(H'\Sigma H\bigr)_{11},
\]
for all $i=1,\ldots,d$. Consequently,
\begin{equation}
\label{eq_zonal_length1}
\int_{\Or(d)} \bigl[\bigl(H'\Sigma H\bigr)_{ii}\bigr]^{p \kappa_i} \dd H = \int_{\Or(d)} \bigl[\bigl(H'\Sigma H\bigr)_{11}\bigr]^{p \kappa_i} \dd H = \frac{C_{(p \kappa_i)}(\Sigma)}{C_{(p \kappa_i)}(I_d)},
\end{equation}
the second equality following from \eqref{eq_zonal}. Applying \eqref{eq_zonal_length1} to \eqref{eq_zonal_Holder}, we obtain
\begin{equation}
\label{eq_zonal_ineq2}
C_\kappa(\Sigma) \le C_\kappa(I_d) \Bigg(\prod_{i=1}^p \frac{C_{(p \kappa_i)}(\Sigma)}{C_{(p \kappa_i)}(I_d)}\Bigg)^{1/p}.
\end{equation}

By \cite[equation~(A.10)]{BagyanRichards},
\[
C_{(p \kappa_i)}(\Sigma) \le \frac{(d^{1/2}/2)_{p \kappa_i}}{(1/2)_{p \kappa_i}} \|\Sigma\|^{p \kappa_i};
\]
also, by \eqref{eq_zonal_identity_1part},
\[
C_{(p \kappa_i)}(I_d) = \frac{(d/2)_{p \kappa_i}}{(1/2)_{p \kappa_i}}.
\]
Substituting the two latter results into \eqref{eq_zonal_ineq2}, we obtain
\begin{align*}
C_\kappa(\Sigma) &\le C_\kappa(I_d) \Bigg(\prod_{i=1}^p \frac{\bigl(d^{1/2}/2\bigr)_{p \kappa_i}}{(1/2)_{p \kappa_i}} \|\Sigma\|^{p \kappa_i} \cdot \frac{(1/2)_{p \kappa_i}}{(d/2)_{p \kappa_i}}\Bigg)^{1/p} \\
&= C_\kappa(I_d) \|\Sigma\|^{|\kappa|} \Bigg(\prod_{i=1}^p \frac{\bigl(d^{1/2}/2\bigr)_{p \kappa_i}}{(d/2)_{p \kappa_i}}\Bigg)^{1/p},
\end{align*}
and this establishes \eqref{eq_zonal_bound}.
\end{proof}

\setcounter{subsection}{2}\setcounter{equation}{32}
% Original source lines 713--828
\subsection[The proofs for Section 2]{The proofs for Section \ref{sec_matrix_bingham}}
\label{sec_proofs_Phi_bounds}

\begin{proof}[Proof of Theorem \ref{thm_Phi_m_bound_Bingham}]
We begin by applying to \eqref{eq_remainder_term} the triangle inequality, the upper bound from Proposition \ref{prop_zonal_bound}, and the inequality $|C_\kappa(A)| \le C_\kappa(A_+)$. Then we obtain
\begin{align}
\label{eq_Phi_m_bound1}
|\Phi_{d,p;m}(A,\Sigma)| &\le \sum_{k=m}^\infty \frac{1}{k!} \sum_{|\kappa|=k} \frac{|C_\kappa(A)|}{C_\kappa(I_d)} \Bigg(\prod_{i=1}^p \frac{\bigl(d^{1/2}/2\bigr)_{p \kappa_i}}{(d/2)_{p \kappa_i}}\Bigg)^{1/p} C_\kappa(I_d) \|\Sigma\|^{|\kappa|} \nonumber \\
&\le \sum_{k=m}^\infty \frac{\|\Sigma\|^k}{k!} \sum_{|\kappa|=k} C_\kappa(A_+) \Bigg(\prod_{i=1}^p \frac{\bigl(d^{1/2}/2\bigr)_{p \kappa_i}}{(d/2)_{p \kappa_i}}\Bigg)^{1/p}.
\end{align}
By \eqref{eq_d_kappa_ratio_ineq3},
\[
\left(\prod_{i=1}^p \frac{\bigl(d^{1/2}/2\bigr)_{p \kappa_i}}{(d/2)_{p \kappa_i}}\right)^{1/p} \le \alpha_p (|\kappa|!)^{1/2} \gamma_1^{|\kappa|} p^{|\kappa|/2} d^{-|\kappa|/2},
\]
and by inserting this bound into \eqref{eq_Phi_m_bound1}, we obtain
\begin{align}
\label{eq_Phi_m_bound2}
\Phi_{d,p;m}(A,\Sigma)| &\le \sum_{k=m}^\infty \frac{\|\Sigma\|^k}{k!} \sum_{|\kappa|=k} C_\kappa(A_+) \alpha_p (|\kappa|!)^{1/2} \gamma_1^{|\kappa|} p^{|\kappa|/2} d^{-|\kappa|/2} \nonumber \\
&= \alpha_p \sum_{k=m}^\infty \frac{\gamma_1^k p^{k/2} d^{-k/2} \|\Sigma\|^k}{(k!)^{1/2}} \sum_{|\kappa|=k} C_\kappa(A_+).
\end{align}

By \eqref{eq_trace_zonal},
$
\sum_{|\kappa|=k} C_\kappa(A_+) = (\tr A_+)^k$,
and, by assumption, $\|\Sigma\| \le \gamma_0 d^{r/2}$. Therefore, we obtain from \eqref{eq_Phi_m_bound2} the inequality
\begin{equation}
\label{eq_Phi_m_bound3}
\Phi_{d,p;m}(A,\Sigma)| \le \alpha_p R_m\bigl(\gamma_0 \gamma_1 p^{1/2} (\tr A_+) d^{-(1-r)/2}\bigr),
\end{equation}
where
\begin{equation}
\label{eq_R_m}
R_m(t) = \sum_{k=m}^\infty \frac{t^k}{(k!)^{1/2}},\qquad t > 0.
\end{equation}
By applying the ratio test, we find that the series \eqref{eq_R_m} converges for all $t$, and therefore~\eqref{eq_Phi_m_bound3} converges for all $d$. This establishes the bound \eqref{eq_Phi_m_bound}.

To obtain a closed-form upper bound for $R_m(t)$ we begin by applying Stirling's well-known inequality for the factorial function \cite[Theorem 1]{Impens}; viz., for all $k \ge 1$,
\[
k! \ge (2\pi)^{1/2} k^{(2k+1)/2} {\rm e}^{-k}.
\]
Inverting Stirling's inequality and applying the result to \eqref{eq_R_m}, we obtain
\begin{equation}
\label{eq_R_m_bound}
R_m(t) \le (2\pi)^{-1/4} \sum_{k=m}^\infty k^{-(2k+1)/4} {\rm e}^{k/2} t^k.
\end{equation}
We now apply the telescoping method to obtain an upper bound for the right-hand side of~\eqref{eq_R_m_bound}. For $j \ge m$, let
\begin{equation}
\label{eq_aj}
a_j = j^{-(2j+1)/4} {\rm e}^{j/2} t^j,
\end{equation}
then we obtain through straightforward algebraic manipulations
\begin{equation}
\label{eq_aj_ratio}
\frac{a_{j+1}}{a_j} = \bigl(1+j^{-1}\bigr)^{-j/2} \frac{j^{1/4}}{(j+1)^{3/4}} {\rm e}^{1/2} t.
\end{equation}
Since the function \smash{$j \mapsto \bigl(1+j^{-1}\bigr)^{-j/2}$}, $j \ge m$, is decreasing, then
\begin{equation}
\label{eq_c_m}
\bigl(1+j^{-1}\bigr)^{-j/2} \le \bigl(1+m^{-1}\bigr)^{-m/2},
\end{equation}
for all $j \ge m$, and by applying the latter inequality to \eqref{eq_aj_ratio}, we obtain
\[
\frac{a_{j+1}}{a_j} \le \bigl(1+m^{-1}\bigr)^{-m/2} {\rm e}^{1/2} t \frac{j^{1/4}}{(j+1)^{3/4}}
= c_m^{1/2} t \frac{j^{1/4}}{(j+1)^{3/4}},
\]
where \smash{$c_m = \bigl(1+m^{-1}\bigr)^{-m} {\rm e}$}. It follows that, for all $k \ge m$,
\begin{align}
\label{eq_a_k}
a_k = a_m \prod_{j=m}^{k-1} \frac{a_{j+1}}{a_j}
\le a_m \prod_{j=m}^{k-1} c_m^{1/2} t \frac{j^{1/4}}{(j+1)^{3/4}}= a_m c_m^{(k-m)/2} t^{k-m} \prod_{j=m}^{k-1} \frac{j^{1/4}}{(j+1)^{3/4}}.
\end{align}
Since \smash{$\prod_{j=m}^{k-1} j = (k-1)!/(m-1)!$}, then
\begin{align}
\label{eq_a_k_factor}
\prod_{j=m}^{k-1} \frac{j^{1/4}}{(j+1)^{3/4}} &= \frac{[(k-1)!/(m-1)!]^{1/4}}{[k!/m!]^{3/4}}
= \frac{m^{3/4} [(m-1)!]^{1/2}}{k^{3/4} [(k-1)!]^{1/2}}.
\end{align}
On applying to \eqref{eq_a_k} the identity \eqref{eq_a_k_factor} and the explicit expression for $a_m$, as given in \eqref{eq_aj}, we obtain
\begin{align}
\label{eq_a_k_bound}
a_k &\le m^{-(2m+1)/4} {\rm e}^{m/2} t^m \cdot c_m^{(k-m)/2} t^{k-m} \frac{m^{3/4} [(m-1)!]^{1/2}}{k^{3/4} [(k-1)!]^{1/2}} \nonumber \\
&= ({\rm e}/m c_m)^{m/2} (m!)^{1/2} \frac{c_m^{k/2} t^k}{k^{3/4} [(k-1)!]^{1/2}}.
\end{align}
Therefore, by \eqref{eq_R_m_bound}, \eqref{eq_aj}, and \eqref{eq_a_k_bound},
\begin{align}
\label{eq_sums_2_bound}
R_m(t) &\le (2\pi)^{-1/4} ({\rm e}/m c_m)^{m/2} (m!)^{1/2} \sum_{k=m}^\infty \frac{c_m^{k/2} t^k}{k^{3/4} [(k-1)!]^{1/2}} \nonumber \\
&\le (2\pi)^{-1/4} ({\rm e}/m c_m)^{m/2} (m!)^{1/2} \Bigg(\sum_{k=m}^\infty \frac{1}{k^{3/2}}\Bigg)^{1/2} \Bigg(\sum_{k=m}^\infty \frac{c_m^k t^{2k}}{(k-1)!}\Bigg)^{1/2},
\end{align}
where the latter inequality follows by the Cauchy--Schwarz inequality.

To obtain an upper bound for the first sum in \eqref{eq_sums_2_bound}, we observe from graphical considerations using Riemann sums that, for $m \ge 2$,
\begin{equation}
\label{eq_zeta_32}
\sum_{k=m}^\infty \frac{1}{k^{3/2}} \le \int_m^\infty \frac{\dd t}{(t-1)^{3/2}} = \frac{2}{(m-1)^{1/2}} \le \left(\frac{8}{m}\right)^{1/2}.
\end{equation}
As for the second sum in \eqref{eq_sums_2_bound}, we apply the inequality
$
\frac{1}{(k+m-1)!} \le \frac{1}{k! (m-1)!}
$
to obtain
\begin{align}
\label{eq_partial_exp}
\sum_{k=m}^\infty \frac{t^{2k}}{(k-1)!} &= \sum_{k=0}^\infty \frac{t^{2(k+m)}}{(k+m-1)!} \le \frac{t^{2m}}{(m-1)!} \sum_{k=0}^\infty \frac{t^{2k}}{k!}
= \frac{t^{2m}}{(m-1)!} \exp\bigl(t^2\bigr).
\end{align}

Applying \eqref{eq_zeta_32} and \eqref{eq_partial_exp} to \eqref{eq_sums_2_bound}, we obtain
\begin{align}
R_m(t) &\le (2\pi)^{-1/4} ({\rm e}/m c_m)^{m/2} (m!)^{1/2} \left(\frac{8}{m}\right)^{1/4} \left(\frac{c_m^m t^{2m}}{(m-1)!} \exp\bigl(c_m t^2\bigr)\right)^{1/2} \nonumber \\
&= (4 {\rm e}/\pi)^{1/4} ({\rm e}/m)^{(m/2)-(1/4)} t^m \exp\bigl(c_m t^2/2\bigr),\label{eq_R_m_bound2}
\end{align}
and by inserting this bound at \eqref{eq_Phi_m_bound3}, we obtain \eqref{eq_Phi_m_boundw}.

Finally, it follows from \eqref{eq_Phi_m_bound} that \smash{$\Phi_{d,p;m}(A,\Sigma) = O\bigl(d^{-(1-r)m/2}\bigr)$} as $d \to \infty$.
\end{proof}
\begin{thebibliography}{99}
\bibitem[1]{BagyanRichards}
Bagyan A., Richards D., Complete asymptotic expansions and the high-dimensional
 {B}ingham distributions, \href{https://doi.org/10.1007/s11749-023-00910-w}{\textit{TEST}} \textbf{33} (2024), 540--563,
 \href{https://arxiv.org/abs/2303.04575}{arXiv:2303.04575}.

\bibitem[2]{Bingham}
Bingham C., An antipodally symmetric distribution on the sphere, \href{https://doi.org/10.1214/aos/1176342874}{\textit{Ann.
 Statist.}} \textbf{2} (1974), 1201--1225.

\bibitem[3]{BinghamCR}
Bingham C., Chang T., Richards D., Approximating the matrix {F}isher and
 {B}ingham distributions: applications to spherical regression and
 {P}rocrustes analysis, \href{https://doi.org/10.1016/0047-259X(92)90072-N}{\textit{J.~Multivariate Anal.}} \textbf{41} (1992),
 314--337.

\bibitem[4]{Chikuse}
Chikuse Y., Statistics on special manifolds, \textit{Lect. Notes Stat.}, Vol.
 174, \href{https://doi.org/10.1007/978-0-387-21540-2}{Springer}, New York, 2003.

\bibitem[5]{DeWaal}
de~Waal D.J., On the normalizing constant for the
 {B}ingham--von~{M}ises--{F}isher matrix distribution, \href{https://doi.org/10520/AJA0038271X_687}{\textit{South African
 Statist.~J.}} \textbf{13} (1979), 103--112.

\bibitem[7]{Dryden}
Dryden I.L., Statistical analysis on high-dimensional spheres and shape spaces,
 \href{https://doi.org/10.1214/009053605000000264}{\textit{Ann. Statist.}} \textbf{33} (2005), 1643--1665, \href{https://arxiv.org/abs/math/0508279}{arXiv:math/0508279}.

\bibitem[8]{Edelman}
Edelman A., Kostlan E., Shub M., How many eigenvalues of a random matrix are
 real?, \href{https://doi.org/10.2307/2152729}{\textit{J.~Amer. Math. Soc.}} \textbf{7} (1994), 247--267.

\bibitem[11]{GrossRichards}
Gross K.I., Richards D.St.P., Special functions of matrix argument.~{I}.
 {A}lgebraic induction, zonal polynomials, and hypergeometric functions,
 \href{https://doi.org/10.2307/2000670}{\textit{Trans. Amer. Math. Soc.}} \textbf{301} (1987), 781--811.

\bibitem[14]{Horn}
Horn R.A., Johnson C.R., Matrix analysis, 2nd ed., \href{https://doi.org/10.1017/CBO9781139020411}{Cambridge University Press},
 Cambridge, 2012.

\bibitem[15]{Impens}
Impens C., Stirling's series made easy, \href{https://doi.org/10.2307/3647856}{\textit{Amer. Math. Monthly}}
 \textbf{110} (2003), 730--735.

\bibitem[16]{James64}
James A.T., Distributions of matrix variates and latent roots derived from
 normal samples, \href{https://doi.org/10.1214/aoms/1177703550}{\textit{Ann. Math. Statist.}} \textbf{35} (1964), 475--501.

\bibitem[17]{James68}
James A.T., Calculation of zonal polynomial coefficients by use of the
 {L}aplace--{B}eltrami operator, \href{https://doi.org/10.1214/aoms/1177698153}{\textit{Ann. Math. Statist.}} \textbf{39}
 (1968), 1711--1718.

\bibitem[18]{James73}
James A.T., The variance information manifold and the functions on it, in
 Multivariate {A}nalysis~{III}, \href{https://doi.org/10.1016/B978-0-12-426653-7.50016-8}{Academic Press}, New York, 1973, 157--169.

\bibitem[20]{JiuKoutschan}
Jiu L., Koutschan C., Calculation and properties of zonal polynomials,
 \href{https://doi.org/10.1007/s11786-020-00458-0}{\textit{Math. Comput. Sci.}} \textbf{14} (2020), 623--640,
 \href{https://arxiv.org/abs/2001.11599}{arXiv:2001.11599}.

\bibitem[22]{JuppMardia2}
Jupp P.E., Mardia K.V., A unified view of the theory of directional statistics,
 1975--1988, \href{https://doi.org/10.2307/1403799}{\textit{Int. Stat. Rev.}} \textbf{57} (1989), 261--294.

\bibitem[23]{KhatriMardia}
Khatri C.G., Mardia K.V., The von {M}ises--{F}isher matrix distribution in
 orientation statistics, \href{https://doi.org/10.1111/j.2517-6161.1977.tb01610.x}{\textit{J.~Roy. Statist. Soc. Ser.~B}} \textbf{39}
 (1977), 95--106.

\bibitem[24]{KumePrestonWood}
Kume A., Preston S.P., Wood A.T.A., Saddlepoint approximations for the
 normalizing constant of {F}isher--{B}ingham distributions on products of
 spheres and {S}tiefel manifolds, \href{https://doi.org/10.1093/biomet/ast021}{\textit{Biometrika}} \textbf{100} (2013),
 971--984.

\bibitem[25]{KumeWood}
Kume A., Wood A.T.A., Saddlepoint approximations for the {B}ingham and
 {F}isher--{B}ingham normalising constants, \href{https://doi.org/10.1093/biomet/92.2.465}{\textit{Biometrika}} \textbf{92}
 (2005), 465--476.

\bibitem[26]{Kushner}
Kushner H.B., On the expansion of {$C^*_\rho(V+I)$} as a sum of zonal
 polynomials, \href{https://doi.org/10.1016/0047-259X(85)90096-X}{\textit{J.~Multivariate Anal.}} \textbf{17} (1985), 84--98.

\bibitem[27]{Ledermann}
Ledermann W., Introduction to group characters, \href{https://doi.org/10.1017/CBO9780511565755}{Cambridge University Press},
 Cambridge, 1977.

\bibitem[29]{MardiaJupp}
Mardia K.V., Jupp P.E., Directional statistics, \textit{Wiley Ser. Probab. Stat.}, \href{https://doi.org/10.1002/9780470316979}{John
 Wiley \& Sons}, Chichester, 1999.

\bibitem[30]{Muirhead}
Muirhead R.J., Aspects of multivariate statistical theory, \textit{Wiley Ser. Probab. Math. Stat.}, John Wiley \& Sons, New York, 1982,
 .

\bibitem[32]{ParkhurstJames}
Parkhurst A.M., James A.T., Zonal polynomials of order~1 through~12, in
 Selected {T}ables in {M}athematical {S}tatistics, {V}ol.~{II}, American
 Mathematical Society, Providence, RI, 1974, 199--388.

\bibitem[33]{Prentice}
Prentice M.J., Antipodally symmetric distributions for orientation statistics,
 \href{https://doi.org/10.1016/0378-3758(82)90025-8}{\textit{J.~Statist. Plann. Inference}} \textbf{6} (1982), 205--214.

\bibitem[34]{Richards85}
Richards D.St.P., Applications of invariant differential operators to
 multivariate distribution theory, \href{https://doi.org/10.1137/0145015}{\textit{SIAM~J. Appl. Math.}} \textbf{45}
 (1985), 280--288.

\bibitem[35]{Richards10}
Richards D.St.P., Functions of matrix argument, in N{IST} {H}andbook of
 {M}athematical {F}unctions, Vol.~35, U.S. Department of Commerce, National
 Institute of Standards and Technology, Washington, DC, 2010, 767--774.

\end{thebibliography}
\end{document}
